% !TeX root = ../概率论与数理统计笔记.tex
% 由主文档通过 \input 载入；请编译主文档。
\setcounter{chapter}{10} % 保留原书章号。
\chapter{充分统计量与统计信息}
\label{chap:ps-information}

本章研究能否压缩样本而保留关于参数的信息，以及如何据此改进估计。
先修内容为统计模型、条件分布和条件期望。充分性为主干，完备性与信息下界可在掌握点估计之后回读。

\section{似然函数与因子分解}
\label{sec:ps-likelihood-sufficiency}
定义\emph{似然函数}（likelihood function）：固定观测 $\mathbf x$，将联合质量函数或密度视为参数 $\theta$ 的函数。
定义\emph{充分统计量}（sufficient statistic），从给定统计量后的样本条件分布不再依赖参数解释其含义。
在共同支配测度下陈述因子分解定理，强调参数依赖的支撑限制属于似然的一部分。

\section{充分性的例子与最小充分性}
\label{sec:ps-minimal-sufficiency}
研究伯努利样本总和、泊松样本总和以及正态模型在不同已知参数条件下的充分统计量。
通过参数影响支撑端点的均匀模型，说明只关注指数项或代数表达式会遗漏信息。
\emph{最小充分性}（minimal sufficiency）及其判据列为选读，明确适用模型与零密度情形。

\section{条件化改进、完备性与无偏估计}
\label{sec:ps-rao-blackwell}
在平方可积无偏估计量与充分统计量的条件下，证明\emph{拉奥--布莱克韦尔定理}（Rao--Blackwell theorem）的方差改进结论。
定义\emph{完备性}（completeness），再陈述莱曼--谢费定理（Lehmann--Scheffé theorem），说明完备充分统计量的无偏函数如何给出唯一的最小方差无偏估计。
将唯一性理解为模型下的几乎处处唯一，并区分充分性与完备性。

\section{费舍尔信息与信息下界}
\label{sec:ps-fisher-information}
定义得分函数（score function）与\emph{费舍尔信息}（Fisher information）。
在可微、积分与微分可交换、得分二阶矩有限等正则条件下，建立信息的常用表示与独立样本的可加性。
陈述克拉默--拉奥下界（Cramér--Rao lower bound）的无偏估计版本；用参数相关支撑的例子说明不能越过正则条件。

\section{例题、习题与学习检查}
规划伯努利概率估计的条件化改进、正态均值的信息计算、均匀模型中常规信息下界的适用性分析。
习题要求区分“统计量包含信息”与“估计值接近真值”，并检查改进结论比较的是哪一类估计量。
完整应用将在第\ref{chap:ps-point-estimation}章与估计方法结合。






